% Generated draft: requires mathematical and source review.
\documentclass[11pt]{article}
\usepackage[margin=25mm]{geometry}
\usepackage{fontspec}
\setmainfont{Latin Modern Roman}
\usepackage{amsmath,amssymb,mathtools,unicode-math}
\setmathfont{Latin Modern Math}
\usepackage{enumitem,needspace,xurl,hyperref}
\setlength{\parindent}{0pt}
\setlength{\parskip}{5pt}
\setlength{\emergencystretch}{3em}
\newcommand{\N}{\mathbb{N}}
\newcommand{\Z}{\mathbb{Z}}
\newcommand{\Q}{\mathbb{Q}}
\newcommand{\R}{\mathbb{R}}
\newcommand{\C}{\mathbb{C}}
\newcommand{\F}{\mathbb{F}}
\newcommand{\A}{\mathbb{A}}
\newcommand{\PP}{\mathbb{P}}
\newcommand{\E}{\mathbb{E}}
\newcommand{\eps}{\varepsilon}
\newcommand{\dd}{\,\mathrm{d}}
\begin{document}
\section*{UnsolvedMath ID: 30004650}
\textbf{Source problem code:} catalogue entry
\subsection{Definitions and mathematical statement}
The supplied record calls this the birational section conjecture, but does not specify the class of base fields \(k\). Accordingly, the formulation below is conditional on the intended choice of that class; it does not assert the conjecture over arbitrary fields.

Under the conventional interpretation of a proper hyperbolic curve, let \(X/k\) be a smooth, proper, geometrically connected curve of genus at least \(2\), and put \(K=k(X)\). These geometric hypotheses spell out a conventional interpretation of the source terminology; they are not separately stated in the supplied record. Choose compatible separable closures \(k^{\mathrm{sep}}\subseteq K^{\mathrm{sep}}\), and write
\[
G_k=\operatorname{Gal}(k^{\mathrm{sep}}/k),\qquad G_K=\operatorname{Gal}(K^{\mathrm{sep}}/K).
\]
Restriction gives a continuous surjection \(\pi:G_K\to G_k\). A continuous section is a continuous group homomorphism \(s:G_k\to G_K\) such that \(\pi\circ s=\operatorname{id}_{G_k}\).

For \(x\in X(k)\), let \(v_x\) be the discrete valuation of \(K\) defined by the local ring \(\mathcal O_{X,x}\). If \(\widetilde v_x\) is a prolongation of \(v_x\) to \(K^{\mathrm{sep}}\), let
\[
D_{\widetilde v_x}=\{g\in G_K:g\text{ preserves }\widetilde v_x\}
\]
be its decomposition subgroup. Interpret a section being induced by \(x\) to mean that its image is contained in such a decomposition subgroup.

With these conventional interpretations, the supplied question is whether
\[
\forall\,s:G_k\to G_K\text{ continuous with }\pi\circ s=\operatorname{id}_{G_k},\quad
\exists!\,x\in X(k)\quad
\exists\,\widetilde v_x\mid v_x\quad
s(G_k)\subseteq D_{\widetilde v_x}.
\]
The uniqueness concerns the rational point attached to a section; it does not assert that a rational point determines a unique literal homomorphism \(G_k\to G_K\). The displayed property is invariant under conjugation of sections by \(\ker\pi\), so it also makes sense for the usual conjugacy classes of sections. The supplied record does not explicitly specify continuity, the conjugacy convention, or the meaning of ``induced.'' Most substantially, the missing restriction on \(k\) prevents an unconditional, fully scoped conjecture statement.
\subsection{Short English statement}
For the intended, unspecified class of base fields, does every continuous section of the absolute Galois group map \(G_{k(X)}\to G_k\) have image in a decomposition subgroup belonging to exactly one \(k\)-rational point of the proper hyperbolic curve \(X\)? The field hypotheses and section conventions require review.
\subsection{Sources}
\begin{itemize}
\item Supplied local UnsolvedMath registry record: permanent ID \(30004650\), title \emph{Birational Sections and Rational Points of Hyperbolic Curves}, exact source code \texttt{catalogue entry}, in the Oberwolfach Reports Open Problems collection. No individual proposer is supplied; the proposed year is recorded as \(2021\).
\item Supplied report citation: \emph{Homotopic and Geometric Galois Theory (2022)}, Oberwolfach Reports, pp.\ 663--744, DOI \texttt{10.4171/owr/2021/12}: \url{https://doi.org/10.4171/owr/2021/12}. The background describes Oberwolfach Reports as documenting mathematical workshops and their open-problem discussions.
\item Supplied literature reference: \emph{Saïdi--Tyler birational section conjecture}, \url{https://doi.org/10.2140/ant.2021.15.435}.
\item Supplied literature reference: \emph{Birational section conjecture with strong assumptions}, \url{https://doi.org/10.1007/s00222-023-01220-6}.
\item Supplied literature reference: \emph{Oberwolfach Report 12/2021}, \url{https://doi.org/10.4171/owr/2021/12}.
\item Literature assessment dated exactly \texttt{2026-08-22}: the supplied curation labels the problem ``partially solved'' and states that Koenigsmann proves the \(p\)-adic birational case, that Saïdi--Tyler reduce finitely generated fields over \(\Q\) to number fields, and that the number-field and hence general birational conjecture remains open. This is supplied dated curation derived from the OWR material, not an independently verified assessment. The source itself requests independent verification of time-sensitive status and cited claims.
\end{itemize}
\end{document}
